\subsubsection{Computational Performance: Speed-Accuracy Frontier}
\label{sec:perf}

In this measurement, we observe the final position deviation of the bodies predicted
by \texttt{SIMON}, in comparison with \texttt{ias15}, over a 100 year time horizon.
The x-axis is the throughput, which is how often the current state of the system is
sampled, \texttt{Rebound}, and the y-axis is RMS position error, on a logarithmic
scale. The total number of samples is fixed across all simulations. The timestep $dt$
is the timestep size the hybrid model uses internally. Points toward the lower-right
of the graph are better (higher throughput, lower error). Throughput is computed as
$\frac{n_{\mathrm{samples}}}{t_{\mathrm{total}}}$.

\begin{figure}[H]
    \centering
    \includegraphics[width=0.82\linewidth]{speed_accuracy_frontier_T100_after_adaptive.png}
    \caption{Speed--accuracy trade-off across timesteps $dt$ for \texttt{SIMON} with
    adaptive sub-stepping. Points towards the lower-right are better (higher
    throughput, lower error). \texttt{SIMON} outperforms \texttt{ias15} at
    $dt = 0.04$ ($1.12\times$ speedup) and $dt = 0.08$ ($2.18\times$ speedup).
    The lowest final RMS error is $0.847$\,AU at $dt = 0.01$.}
    \label{fig:frontier}
\end{figure}

The \texttt{ias15} baseline completes in $0.080$ seconds. \texttt{SIMON} outperforms
\texttt{ias15} at $dt = 0.04$ ($1.12\times$ speedup) and $dt = 0.08$ ($2.18\times$
speedup). The lowest final RMS error is $0.847$ at $dt = 0.01$. This is the optimal
point for accuracy, but it is approximately $3.4\times$ slower than \texttt{ias15}.
The errors at $dt = 0.005$, $0.02$, and $0.08$ are all approximately $1.02$,
reflecting the chaotic saturation scale. After 100 years, any two different integrations
of the same system will diverge to roughly one orbital radius, regardless of timestep.
$dt = 0.04$ shows a slightly higher error of $1.54$, which is a chaotic fluctuation
(a different phase of the orbit was reached at $T = 100$).

Table~\ref{tab:frontier} provides the full quantitative results for the frontier sweep.

\begin{table}[H]
\centering
\caption{Speed-accuracy frontier results for \texttt{SIMON} with adaptive
sub-stepping, $T = 100$\,yr. \texttt{ias15} baseline: $0.080$\,s total,
$\lambda = 0.1655$\,yr$^{-1}$ at $dt = 0.04$\,yr. Bold row = operational timestep.}
\label{tab:frontier}
\resizebox{\columnwidth}{!}{%
\begin{tabular}{@{}lllllll@{}}
\toprule
$dt$ (yr) & Final RMS (AU) & Throughput (samp/s) & Time/step ($\mu$s) & Total time (s) & Steps & Speedup \\
\midrule
$0.005$  & $1.022$ & $9{,}285$   & $26.9$ & $0.538$ & $20{,}000$ & $0.15\times$ \\
$0.010$  & $0.847$ & $18{,}307$  & $27.3$ & $0.273$ & $10{,}000$ & $0.29\times$ \\
$0.020$  & $1.022$ & $35{,}672$  & $28.0$ & $0.140$ & $5{,}000$  & $0.57\times$ \\
$\mathbf{0.040}$ & $\mathbf{1.543}$ & $\mathbf{69{,}880}$ & $\mathbf{28.6}$ & $\mathbf{0.072}$ & $\mathbf{2{,}500}$ & $\mathbf{1.12\times}$ \\
$0.080$  & $1.022$ & $135{,}629$ & $29.5$ & $0.037$ & $1{,}250$  & $2.18\times$ \\
\bottomrule
\end{tabular}%
}
\end{table}

The per-step time ranges from $26.9\,\mu$s at $dt = 0.005$ to $29.5\,\mu$s at
$dt = 0.08$: a range of only $2.6\,\mu$s across a $16\times$ change in timestep.
This nearly flat profile indicates that the cost of an individual \texttt{SIMON}
step is approximately constant and is dominated by force evaluation, pairwise
distance checks, and occasional adaptive sub-stepping. Therefore, the observed
runtime improvement at larger $dt$ comes primarily from taking fewer macro steps,
not from making each individual step substantially cheaper. For a fixed
$T = 100$\,yr horizon, the number of macro steps scales approximately as $T/dt$,
so total runtime scales roughly as $(T/dt)$ times the near-constant per-step cost.
